\subsection{局部紧交换群的酉特征标}

定义 1.　G 的酉特征标是从 G 到模为 1 的复数所成的乘法群 U 中的一个连续同态。

换句话说，酉特征标是 G 上的一个连续复值函数 \(\chi\)，满足：

\[
\chi (x y) = \chi (x) \chi (y), \qquad | \chi (x) | = 1 \qquad (x, y \in \mathrm{G}).
\]

在本章中，我们常把"酉特征标"简称为"特征标"。

设 E 为一维 Hilbert 空间，\(\chi\) 为 G 的一个酉特征标。把 \(x \in \mathrm{G}\) 映到 E 中比为 \(\chi(x)\) 的位似映射的映射是 G 在 E 上的一个连续线性等距表示。反之，G 在 E 上的每个有界连续线性表示都可由此过程得到，特别地它是酉的。

两个酉特征标的乘积、一个酉特征标的逆以及等于 1 的常函数都是酉特征标，这是直接的。因此，G 的酉特征标所成的集合 \(\widehat{\mathbf{G}}\) 对乘法构成一个群。这个群是交换的。另一方面，映射 \((\chi_1,\chi_2)\mapsto \chi_1\chi_2^{-1} = \chi_1\overline{\chi}_2\) 对紧开拓扑是连续的，从而赋有紧收敛拓扑的 \(\widehat{\mathbf{G}}\) 是一个拓扑群（TG，X，第 6 页，推论 2 与注 1）。

定义 2.　拓扑群 \(\widehat{\mathbf{G}}\) 称为 G 的对偶群。

由于 G 是局部紧的，映射 \((x,\chi)\mapsto\chi(x)\) 在 \(G\times\widehat{G}\) 上连续（TG，X，第 28 页，定理 3）。

回忆 \(\mathcal{M}^{1}(G)\) 表示 G 上有界复测度所成的幺对合 Banach 代数（I，第 99 页，例 4）。对每个有界复测度 \(\mu\in\mathcal{M}^{1}(G)\) 与每个 \(\chi\in\widehat{G}\)，我们记

\[
\chi (\mu) = \int_ {\mathrm{G}} \chi (x) d \mu (x)\tag{1}
\]

（参见 INT，VIII，§2，第 6 节）。

引理 1.　对每个 \(\chi \in \widehat{\mathbf{G}}\)，映射 \(\mu \mapsto \chi(\mu)\) 是对合 Banach 代数 \(\mathcal{M}^1(\mathbf{G})\) 的一个 Hermite 特征标。

由 INT，VIII，§3，第 3 节，命题 11，映射 \(\mu \mapsto \chi(\mu)\) 是对合 Banach 代数 \(\mathcal{M}^{1}(G)\) 的一个特征标。此外，我们有：

\[
\chi (\mu^ {*}) = \int_ {\mathrm{G}} \chi (x ^ {- 1}) d \overline {{\mu}} (x) = \int_ {\mathrm{G}} \overline {{\chi (x)}} d \overline {{\mu}} (x) = \overline {{\chi (\mu)}}
\]

因此这个特征标是 Hermite 的。

这样我们定义了一个从 \(\widehat{G}\) 到 \(\mathsf{X}(\mathcal{M}^{1}(\mathrm{G}))\) 中的映射；它将称为典范映射。

设 \(\chi \in \widehat{\mathbf{G}}\)。\(\mu \mapsto \chi(\mu)\) 在 \(\mathrm{L}^1(\mathrm{G})\) 上的限制非零（参见 INT，VIII，§2，第 7 节，命题 10）。因此，通过限制到对合 Banach 子代数 \(\mathrm{L}^1(\mathrm{G})\)，我们得到一个从 \(\widehat{\mathbf{G}}\) 到 \(\mathsf{X}(\mathrm{L}^1(\mathrm{G}))\) 中的映射，称为典范映射。它把 \(\chi \in \widehat{\mathbf{G}}\) 对应到 Hermite 特征标

\[
f \mapsto \chi (f) = \chi (f \cdot d x) = \int_ {\mathrm{G}} f (x) \chi (x) d x \qquad (f \in \mathrm{L} ^ {1} (\mathrm{G}))\tag{2}
\]

即 \(L^{1}(G)\) 的特征标。

命题 1.　从 \(\widehat{\mathbf{G}}\) 到 \(\mathsf{X}(\mathrm{L}^{1}(\mathbf{G}))\) 的典范映射是一个同胚。

把这个典范映射记作 ev，并对 \(\chi \in \widehat{\mathbf{G}}\)，用 \(\mathrm{ev}_{\chi}\) 表示特征标 \(\chi\) 在 ev 下的像，即 Hermite 特征标 \(f \mapsto \chi(f)\)，它定义在 \(\mathrm{L}^1(\mathrm{G})\) 上。把 ev 视为从 \(\widehat{\mathbf{G}}\) 到 \(\mathrm{L}^1(\mathrm{G})\) 的对偶中的映射时，它是 \(\widehat{\mathbf{G}}\) 到 \(\mathrm{L}^{\infty}(\mathrm{G})\)（赋有弱拓扑 \(\sigma(\mathrm{L}^{\infty}(\mathrm{G}), \mathrm{L}^{1}(\mathrm{G}))\)）中的单射与 \(\mathrm{L}^{\infty}(\mathrm{G})\) 到赋有点态收敛拓扑的 \(\mathrm{L}^1(\mathrm{G})\) 的对偶中的单射的复合。由于 \(\widehat{\mathbf{G}}\) 是 \(\mathrm{L}^{\infty}(\mathrm{G})\) 的有界子集，由 Lebesgue 定理（INT，IV，§3，第 7 节，定理 6），第一个映射连续。第二个映射按定义也是连续的。这就证明了映射 ev 连续。

若 \(\chi \in \widehat{\mathrm{G}}\) 且 \(f \in \mathrm{L}^1(\mathrm{G})\)，则有 \(\mathrm{ev}_{\chi}(\varepsilon_x * f) = \chi(x)\mathrm{ev}_{\chi}(f)\)。取 \(f\) 使 \(\mathrm{ev}_{\chi}(f) \neq 0\)，便推出映射 ev 是单射。

设 \(\zeta\in\mathsf{X}(\mathrm{L}^{1}(\mathrm{G}))\)，并设 \(f\in\mathrm{L}^{1}(\mathrm{G})\) 满足 \(\zeta(f)\neq0\)。我们用如下方式定义映射 \(\chi\colon\mathrm{G}\to\mathbf{C}\)：对 \(x\in\mathrm{G}\)，令

\[
\chi (x) = \frac {\zeta (\varepsilon_ {x} * f)}{\zeta (f)}.\tag{3}
\]

来定义它。我们有 \(\chi(e) = 1\)。由于映射 \(x \mapsto \varepsilon_x * f = \gamma(x)(f)\)（从 G 到 L \(^1\) (G)）连续（INT，VIII，§2，第 5 节，命题 8），映射 \(\chi\) 连续。它是有界的，因为对所有 \(x \in \mathbf{G}\) 有

\[
| \chi (x) | \leqslant \frac {\| \varepsilon_ {x} * f \|}{| \zeta (f) |} = \frac {\| f \|}{| \zeta (f) |}
\]

（I，第 29 页，定理 1 及 INT，VIII，同前）。

现设 \(\mathfrak{B}\) 为由紧邻域组成的 \(e\) 的邻域滤子的一个基。对每个 \(V \in \mathfrak{B}\)，设 \(g_V\) 为一个在 \(V\) 外为零且积分等于 1 的正连续函数（II，第 200 页，引理 1）。于是对每个函数 \(h \in L^1(G)\) 有

\[
\varepsilon_ {x} * h = \lim _ {\mathrm{V}, \mathfrak {B}} (\varepsilon_ {x} * h) * g _ {\mathrm{V}} = \lim _ {\mathrm{V}, \mathfrak {B}} (\varepsilon_ {x} * g _ {\mathrm{V}} * h),
\]

在 L \(^{1}\) (G) 中成立，其中极限是沿 \(\mathfrak{B}\) 的截口滤子取的（INT，VIII，§4，第 7 节，命题 20）。特别地，由于 \(\zeta(\varepsilon_x * g_V * f) = \zeta(\varepsilon_x * g_V)\zeta(f)\)，由此推出

\[
\chi (x) = \lim _ {\mathrm{V}, \mathfrak {B}} \zeta (\varepsilon_ {x} * g _ {\mathrm{V}}).
\]

对每个 \(h\in \mathrm{L}^1 (\mathrm{G})\)，我们得到

\[
\zeta (\varepsilon_ {x} * h) = \lim _ {\mathrm{V}, \mathfrak {B}} \zeta (\varepsilon_ {x} * g _ {\mathrm{V}} * h) = \zeta (h) \lim _ {\mathrm{V}, \mathfrak {B}} \zeta (\varepsilon_ {x} * g _ {\mathrm{V}}) = \chi (x) \zeta (h).
\]

从而，对 \(x, y \in G\) 有：

\[
\chi (x y) = \frac {\zeta (\varepsilon_ {x} * \varepsilon_ {y} * f)}{\zeta (f)} = \frac {\chi (x) \zeta (\varepsilon_ {y} * f)}{\zeta (f)} = \chi (x) \chi (y),
\]

这就证明了 \(\chi\) 是从 G 到 \(\mathbf{C}^*\) 的一个同态。由于 \(\chi\) 有界且连续，它是 G 的一个酉特征标。此外，若 \(g \in \mathrm{L}^{1}(\mathrm{G})\)，则有

\[
g * f = \int_ {\mathrm{G}} (\varepsilon_ {x} * f) g (x) d x
\]

在 L \(^{1}\) (G) 中成立（INT，VIII，§1，第 5 节，命题 7），由此得

\[
\begin{array}{c} \zeta (g) \zeta (f) = \zeta (g * f) = \int_ {\mathrm{G}} \zeta (\varepsilon_ {x} * f) g (x) d x \\ = \zeta (f) \int_ {\mathrm{G}} \chi (x) g (x) d x = \mathrm{ev} _ {\chi} (g) \zeta (f) \end{array}
\]

（INT，VI，§1，第 1 节，命题 1），这表明 \(\zeta = \mathrm{ev}_{\chi}\)。因此 ev 是满的，从而是双射。

最后证明逆映射 \(ev^{-1}\) 连续。设 \(\zeta \in \mathsf{X}(\mathrm{L}^{1}(\mathrm{G}))\)。设 \(f \in \mathrm{L}^{1}(\mathrm{G})\) 为使 \(\zeta(f) \neq 0\) 的函数。由 \(\xi \in \mathsf{X}(\mathrm{L}^{1}(\mathrm{G}))\)（满足 \(\xi(f) \neq 0\)）所成的集合 W 是 \(\zeta\) 在 \(\mathsf{X}(\mathrm{L}^{1}(\mathrm{G}))\) 中的一个开邻域。对每个 \(\xi \in W\)，前文的论证表明 \(\mathrm{ev}^{-1}(\xi)\) 是特征标

\[
x \mapsto \frac {\xi (\varepsilon_ {x} * f)}{\xi (f)}.
\]

设 \(\mathfrak{F}\) 为 \(W \subset X(L^{1}(G))\) 上收敛于 \(\zeta\) 的一个滤子。由于集合 \(X(L^{1}(G))\) 在 \(L^{\infty}(G)\) 中有界从而等度连续，单收敛的一致结构与紧收敛的一致结构一致（TG，X，第 16 页，定理 1）。设 K 为 G 的一个紧子集。诸 \(\varepsilon_{x} * f\)（\(x \in K\)）所成的集合在 \(L^{1}(G)\) 中是紧的（INT，VIII，§2，第 5 节，命题 8）。于是我们有

\[
\lim _ {\xi , \mathfrak {F}} \xi (\varepsilon_ {x} * f) = \zeta (\varepsilon_ {x} * f)
\]

对 \(x \in K\) 一致地成立。由此推出

\[
\lim _ {\xi , \mathfrak {F}} \mathrm{ev} ^ {- 1} (\xi) = \mathrm{ev} ^ {- 1} (\zeta),
\]

即 ev \(^{-1}\) 在 \(\zeta\) 处连续。命题证毕。

从现在起，我们把 G 的酉特征标 \(\chi\) 与特征标 \(f\mapsto\int_{\mathrm{G}}f(x)\chi(x)dx\)（\(\mathrm{L}^{1}(\mathrm{G})\) 上的）等同起来。

注.　1) *命题 1 中从 \(\widehat{G}\) 到 \(X(L^{1}(G))\) 的映射的双射性，是局部紧群 H（不必交换）的连续表示与代数 \(L^{1}(H)\) 的连续表示之间对应关系的一个特殊情形。*

\begin{enumerate}
\def\labelenumi{\arabic{enumi})}
\setcounter{enumi}{1}
\tightlist
\item
  从 \(\widehat{G}\) 到 \(\mathsf{X}(\mathcal{M}^{1}(\mathsf{G}))\) 的典范映射一般不是满射（II，第 308 页，习题 14）。
\end{enumerate}

推论 1.　L \(^{1}\) (G) 的每个特征标都是 Hermite 的。从 X(Stell(G))（I，第 125 页，定义 9）到 X(L \(^{1}\) (G)) 的典范映射是一个同胚。

第一个断言由命题 1 与引理 1 限制到 \(L^{1}(G)\) 得到。第二个断言由第一个断言以及 I，第 124 页，命题 20 的推论得到。

推论 2.　拓扑群 \(\widehat{\mathbf{G}}\) 是局部紧的。

事实上，\(\mathsf{X}(\mathrm{L}^{1}(\mathrm{G}))\) 是局部紧的（I，第 29 页，定理 1 的推论）。

我们将把 \(\widehat{G}\) 与 \(\mathsf{X}(\mathrm{L}^1 (\mathrm{G}))\) 以及 \(\mathsf{X}(\mathrm{Stell}(\mathrm{G}))\) 等同起来。对 \(x\in G\) 与 \(\chi \in \widehat{G}\)，我们用 \(\langle \chi ,x\rangle\) 表示复数 \(\chi (x)\)，它属于 \(\mathbf{U}\)。

我们称 \(x\) 与 \(\chi\) 正交，如果 \(\langle \chi, x \rangle = 1\)。设 A 为 G（相应地，\(\widehat{\mathbf{G}}\)）的一个子集；\(\widehat{\mathbf{G}}\)（相应地，G）中与 A 正交的元素所成的集合是 \(\widehat{\mathbf{G}}\)（相应地，G）的一个闭子群，称为 A 的正交子群，记作 \(\mathrm{A}^{\perp}\)。G 的正交子群仅由 \(e\) 组成。

对 \(x \in G\)，设 \(\eta(x)\) 为从 \(\widehat{G}\) 到 U 的、由 \(\chi \mapsto \langle \chi, x \rangle\) 定义的映射。由 \(\widehat{G}\) 中乘法的定义，映射 \(\eta(x)\) 是一个群同态。由于映射 \((x, \chi) \mapsto \langle \chi, x \rangle\)（从 \(G \times \widehat{G}\) 到 U）连续（TG，X，第 28 页，定理 3），它连续。这样我们定义了一个映射 \(\eta\)，称为典范映射，它从 G 映到二重对偶群 \(\widehat{\widehat{G}}\) 中；它是一个群同态。此外，映射 \(\eta\) 连续（TG，X，第 28 页，定理 3）。我们稍后将证明（II，第 220 页，定理 2）\(\eta\) 是从 G 到 \(\widehat{\widehat{G}}\) 上的拓扑群同构。

设 G 与 H 为局部紧交换群，\(\varphi\colon\mathrm{G}\to\mathrm{H}\) 为一个拓扑群态射。对每个 \(\chi\in\widehat{\mathrm{H}}\)，映射 \(\chi\circ\varphi\) 是 G 的一个特征标，记作 \(\widehat{\varphi}(\chi)\)。这一定义可以写成公式

\[
\langle \chi , \varphi (x) \rangle = \langle \widehat {\varphi} (\chi), x \rangle\tag{4}
\]

对所有 \(\chi\in\widehat{H}\) 与 \(x\in G\) 成立。由此推出 \(\widehat{\varphi}\) 是从拓扑群 \(\widehat{H}\) 到拓扑群 \(\widehat{G}\) 的一个态射；我们称 \(\widehat{\varphi}\) 为态射 \(\varphi\) 的对偶。

设 K 为局部紧交换群，\(\psi\colon\mathrm{H}\to\mathrm{K}\) 为一个拓扑群态射。定义表明 \(\widehat{\psi\circ\varphi}=\widehat{\varphi}\circ\widehat{\psi}\)。若 \(\varphi\) 是 G 的恒等映射，则 \(\widehat{\varphi}\) 是 \(\widehat{\mathrm{G}}\) 的恒等映射。特别地，若 \(\varphi\) 是拓扑群同构，则 \(\widehat{\varphi}\) 也是，且 \(\widehat{\varphi}^{-1}\) 是 \(\varphi^{-1}\) 的对偶。

引理 2.　设 G 与 H 为局部紧交换群，\(f\colon\mathrm{H}\to\mathrm{G}\) 为一个拓扑群态射。\(\widehat{f}\) 的核是 \(f\) 的像的正交子群。

按定义，\(\chi \in \mathrm{Ker}(\widehat{f})\) 当且仅当 \(\chi\) 在 \(f\) 的像上的限制是平凡的。

命题 2.　设 \(n \geqslant 0\) 为整数，\(\mathrm{G}_{1}, \ldots, \mathrm{G}_{n}\) 为局部紧交换群。设 G 为诸群 \(\mathrm{G}_{j}\)（\(1 \leqslant j \leqslant n\)）的乘积群。对 \(1 \leqslant j \leqslant n\)，设 \(\lambda_{j}\) 为 \(\mathrm{G}_{j}\) 到 G 中的单射，它把 \(x \in \mathrm{G}_{j}\) 映到元素 \((x_{k})\)（\(x_{k} = e\) 若 \(k \neq j\)，而 \(x_{j} = x\)）。映射

\[
(\widehat {\lambda} _ {j}) _ {1 \leqslant j \leqslant n} \colon \widehat {\mathrm{G}} \to \prod_ {1 \leqslant j \leqslant n} \widehat {\mathrm{G}} _ {j}
\]

是一个拓扑群同构。

设 \(m\) 为从 \(\widehat{\mathbf{G}}^n\) 到 \(\widehat{\mathbf{G}}\) 的乘积映射，并对每个 \(j\)（\(1 \leqslant j \leqslant n\)），设 \(\pi_j\) 为 \(\mathbf{G}\) 到 \(\mathbf{G}_j\) 上的投影。设 \(\mu\) 为拓扑群态射 \(m \circ (\widehat{\pi}_j)_j\)，它从 \(\prod \widehat{\mathbf{G}}_j\) 到 \(\widehat{\mathbf{G}}\)，于是

\[
\langle \mu ((\chi_ {j})), (x _ {j}) \rangle = \prod_ {j = 1} ^ {n} \langle \chi_ {j}, x _ {j} \rangle .
\]

映射 \(\mu\) 连续，并且可以验证 \(\mu\) 与 \((\widehat{\lambda}_j)\) 互为逆双射。命题得证。

注.　紧交换群的无限乘积的对偶群的计算是下文 II，第 234 页，推论 4 的主题。至于作为局部紧群的任意乘积的局部紧交换群的情形，则可由这两个命题得到，因为在这种乘积中除有限多个因子外都是紧的（TG，I，第 66 页，命题 14，b））。

